\tikzset{
  ctrlBlock/.style={
    draw,
    rounded corners=2pt,
    align=center,
    minimum height=8mm,
    minimum width=24mm,
    inner xsep=4pt,
    inner ysep=3pt,
    font=\small
  },
  ctrlSmall/.style={
    draw,
    rounded corners=2pt,
    align=center,
    minimum height=7mm,
    minimum width=20mm,
    inner xsep=3pt,
    inner ysep=2pt,
    font=\footnotesize
  },
  ctrlSum/.style={
    draw,
    circle,
    inner sep=0pt,
    minimum size=5.5mm,
    font=\small
  },
  ctrlGroup/.style={
    draw,
    dashed,
    rounded corners=3pt,
    inner sep=6pt
  },
  ctrlArrow/.style={
    -{Latex[length=2.2mm]},
    line width=0.45pt
  },
  ctrlDashArrow/.style={
    -{Latex[length=2.2mm]},
    dashed,
    line width=0.45pt
  }
}


\begin{figure}[htbp]
  \centering
  \resizebox{0.98\textwidth}{!}{%
  \begin{tikzpicture}
    \node[ctrlBlock] (cmd) at (0,0) {头车输入\\$v_1,\omega_1$};
    \node[ctrlBlock] (pose) at (3.0,0) {各单体位姿\\$x_i,y_i,\phi_i$};
    \node[ctrlBlock] (joint) at (6.0,0) {铰接角计算\\$\theta_i$};
    \node[ctrlBlock] (constraint) at (9.0,0) {铰接点速度\\一致约束};
    \node[ctrlBlock] (prop) at (12.0,0) {后车期望输入\\$v_{i+1},\omega_{i+1}$};

    \node[ctrlSmall] (leg) at (3.0,-1.45) {腿部运动学\\$q_j \rightarrow L_0,\Phi_0$};
    \node[ctrlSmall] (seq) at (9.0,-1.45) {逐级传播\\形成链式期望序列};

    \draw[ctrlArrow] (cmd) -- (pose);
    \draw[ctrlArrow] (pose) -- (joint);
    \draw[ctrlArrow] (joint) -- (constraint);
    \draw[ctrlArrow] (constraint) -- (prop);
    \draw[ctrlArrow] (pose) -- (leg);
    \draw[ctrlArrow] (prop) |- (seq);
    \draw[ctrlDashArrow] (seq.south) -- ++(0,-0.7) -| node[pos=0.24, below, font=\footnotesize] {更新下一单体} ($(leg.west)+(-0.35,-0.05)$) |- (pose.west);
  \end{tikzpicture}
  }%
  \caption{链式系统运动学建模框图}
  \label{fig:ch2_chain_kinematics}
\end{figure}
